\documentclass[a4paper]{exam}
\usepackage[svgnames]{xcolor}
\usepackage{amsmath,amssymb,amsthm,titlesec,minted}

\newcommand
\lralign[2]{%
  \noindent\makebox[\textwidth]{\(\displaystyle#1\)\hfill\ifx&#2&\else(#2)\fi}\vspace{2ex}
}

\renewcommand{\thesubsection}{\alph{subsection}.}
\renewcommand{\thesubsubsection}{\roman{subsubsection}.}

\definecolor{faintgreen}{rgb}{0.5,1,0.5}
\definecolor{faintred}{rgb}{1,0.5,0.5}
\definecolor{faintblue}{rgb}{0.5,0.5,1}
\definecolor{faintyellow}{rgb}{1,1,0.5}
\definecolor{codebg}{RGB}{40,42,54}

\titleformat{\section}{\normalfont\Large\bfseries}{}{0pt}{}
\titleformat{\subsection}{\normalfont\large\bfseries}{}{0pt}{}
\titleformat{\subsubsection}{\normalfont\normalsize\bfseries\itshape}{}{0pt}{}
\titlespacing*{\section}{1em}{1.6\baselineskip}{0.6\baselineskip}
\titlespacing*{\subsection}{2em}{1.2\baselineskip}{0.4\baselineskip}
\titlespacing*{\subsubsection}{4em}{1.0\baselineskip}{0.3\baselineskip}

\setminted{style=dracula,bgcolor=codebg,breaklines,autogobble,fontsize=\small,tabsize=4,baselinestretch=1}
\NewDocumentCommand{\codefromfile}{O{}mm}{\par\vspace{0.4em}\inputminted[#1]{#2}{#3}\par\vspace{0.4em}}
\NewDocumentEnvironment{Section}{m}
{%
  \section*{#1}%
  \begingroup
  \leftskip=1em
}
{%
  \par
  \endgroup
}

\NewDocumentEnvironment{Subsection}{m}
{%
  \subsection*{#1}%
  \begingroup
  \leftskip=2em
}
{%
  \par
  \endgroup
}

\setlength{\parindent}{0in}
\setlength{\oddsidemargin}{0in}
\setlength{\textwidth}{6.5in}
\setlength{\textheight}{8.8in}
\setlength{\topmargin}{0in}
\setlength{\headheight}{18pt}

\printanswers
\title{Week 4}
\author{Radhesh Goel}
\date{\today}

\begin{document}
\maketitle

\section*{Question 1.}
\begin{solution}
  The combustion reaction is
  \[
    \mathrm{CH}_4 + \alpha\mathrm{O}_2 + 4\alpha\mathrm{N}_2 \rightarrow \mathrm{CO}_2
                                                                         + 2\mathrm{H}_2\mathrm{O}
                                                                         + \beta\mathrm{NO}
                                                                         + \gamma\mathrm{N}_2.
  \]

  Equating the number of oxygen atoms on each side,
  \[
    2\alpha = 2 + 2 + \beta.
  \]
  Therefore,
  \[
    \boxed{2\alpha-\beta=4}.
  \]

  Equating the number of nitrogen atoms on each side,
  \[
    2(4\alpha) = \beta + 2\gamma.
  \]
  Therefore,
  \[
    \boxed{8\alpha-\beta-2\gamma=0}.
  \]
\end{solution}

\section*{Question 2.}
\begin{solution}
  The three equations are
  \[
    \begin{aligned}
      2\alpha - \beta           & = 4, \\
      8\alpha - \beta - 2\gamma & = 0, \\
      2\beta - \gamma           & = 0
    \end{aligned}
  \]

  Let
  \[
    \mathbf{x} = \begin{bmatrix}
                   \alpha \\
                   \beta \\
                   \gamma
                 \end{bmatrix}
  \]

  The equations can be written as
  \[
    \underbrace{
        \begin{bmatrix}
          2 & -1 & 0\\
          8 & -1 & -2\\
          0 & 2 & -1
        \end{bmatrix}
      }_{\mathbf{A}} \underbrace{
        \begin{bmatrix}
          \alpha\\
          \beta\\
          \gamma
        \end{bmatrix}
      }_{\mathbf{x}}
    = \underbrace{
        \begin{bmatrix}
          4\\
          0\\
          0
        \end{bmatrix}
      }_{\mathbf{b}}
  \]

  Therefore,
  \[
    \boxed{
        \mathbf{A}
        =
        \begin{bmatrix}
          2 & -1 & 0\\
          8 & -1 & -2\\
          0 & 2 & -1
        \end{bmatrix},
        \qquad
        \mathbf{b}
        =
        \begin{bmatrix}
          4\\
          0\\
          0
        \end{bmatrix}
      }
  \]
\end{solution}
\end{document}
